% Bagéan: teori-himpunan
% Bab: ordinal
% Subbagéan: pambuka

\documentclass[../../../include/open-logic-section]{subfiles}

\begin{document}

\olfileid{sth}{ordinals}{intro}

\olsection{Pambuka}

Ing \olref[z][]{chap}, kita wis nganggep ana tataran
tanpa wates ing hierarki, arupa \stagesinf{} (delengen
uga Aksioma Tanpa Wates). Nanging, amarga \stagessucc{},
kita ora bisa mandheg ing tataran tanpa wates iku;
kita kudu nerusake. Dadi ing tataran sawisé tataran
tanpa wates kapisan, kita mbentuk kabeh kumpulan
sing bisa kawangun saka himpunan kang wis kasedhiya
ing tataran tanpa wates kapisan; banjur mbaleni maneh;
lan mbaleni maneh; lan mbaleni maneh; \dots

Kanthi ora langsung, ing kene kita wis wiwit nggunakake
gagasan ``intuitif'' bab wilangan, kang ngidini anane
wilangan \emph{sawisé} kabeh wilangan asli. Mligine,
gagasan iku yaiku \emph{ordinal transfinit}. Tujuan
bab iki yaiku njlentrehake gagasan kasebut kanthi
luwih ketat. Kita bakal nliti gagasan umum bab ordinal,
banjur netepake kanthi cetha sawenehing himpunan
minangka ordinal kita.

\end{document}
